\subsection{Metodologia}

% A metodologia adotada neste trabalho é de natureza teórico-computacional, fundamentada no desenvolvimento matemático e na validação por meio de simulação numérica.
% Inicialmente, realizou-se uma revisão bibliográfica sobre dinâmica orbital relativa, robótica espacial e operações de \gls{rvdb}, com o objetivo de identificar lacunas na integração entre essas áreas.

% Na sequência, foi desenvolvida uma formulação matemática integrada, contemplando a dinâmica orbital relativa da espaçonave perseguidora em relação à espaçonave alvo e a dinâmica acoplada entre a espaçonave e seu \gls{mr}.
% A modelagem cinemática do manipulador foi conduzida por meio do método de \gls{dh}, enquanto a dinâmica acoplada foi formulada utilizando abordagem lagrangiana baseada na energia cinética do conjunto espaçonave--\gls{mr}.

% O modelo foi implementado computacionalmente em Matlab e Simulink, incluindo a incorporação de lei de controle do tipo \gls{pd} para o movimento orbital, atitude e movimento das juntas do \gls{mr}.
% Para a aproximação de longa distância foram realizadas simulações numéricas para diferentes condições iniciais, que permitiu avaliar o comportamento dinâmico da espaçonave durante esta fase e como iria interferir nas fases seguintes, como a aproximação de curta distância e aproximação final.

% A avaliação dos resultados foi conduzida por meio da análise dos perfis temporais das variáveis orbitais e de atitude, com ênfase na influência do \gls{mr} sobre a dinâmica da espaçonave perseguidora.
% Em resumo, foi realizado o seguinte:

% \begin{itemize}
%     \item Revisão bibliográfica sobre dinâmica orbital relativa, robótica espacial e operações de \gls{rvdb}.
%     \item Desenvolvimento da formulação matemática integrada.
%     \item Modelagem cinemática do \gls{mr} via \gls{dh}.
%     \item Formulação dinâmica acoplada por abordagem lagrangiana.
%     \item Implementação computacional em Matlab e Simulink.
%     \item Implementação de lei de controle \gls{pd}.
%     \item Realização de simulações numéricas sob diferentes condições iniciais.
%     \item Análise qualitativa dos resultados por inspeção dos perfis temporais das variáveis dinâmicas.
% \end{itemize}

A metodologia adotada é de natureza teórico-computacional, fundamentada no desenvolvimento matemático e na validação por meio de simulação numérica. Inicialmente, realizou-se uma revisão da literatura sobre dinâmica orbital relativa, robótica espacial e operações de RVD/B, com o objetivo de identificar lacunas na integração entre essas áreas.
Na sequência, foi desenvolvida uma formulação matemática integrada, contemplando a dinâmica orbital relativa da perseguidora em relação ao alvo e a dinâmica acoplada entre a espaçonave e seu manipulador. 
A modelagem cinemática do manipulador foi conduzida pela convenção de \gls{dh}, enquanto a dinâmica acoplada foi formulada via abordagem lagrangiana baseada na formulação via energia e coordenadas generalizadas. Com essa metodologia derivou-se as equações movimento acopladas energia através do uso da energia cinética do sistema espaçonave tipo manipulador robótico na configuração de corpo não rígido.
O modelo foi simulado computacionalmente em ambiente Matlab e Simulink, incluindo controlador PD para movimento orbital, atitude e juntas do manipulador. Foram realizadas simulações numéricas em diferentes condições iniciais, permitindo avaliar o comportamento dinâmico da perseguidora nas fases de aproximação de longa distância, curta distância e final.
A avaliação dos resultados foi conduzida por meio da análise dos perfis temporais das variáveis orbitais e de atitude, com ênfase na influência do manipulador sobre a dinâmica da perseguidora.
